\section{Recovering bounded operators from spectral vectors}
\label{sec:recovery}

\Cref{rec:bounded} converts a positive averaged quantity on narrow modular
spectral bands into a positive quantity on uniformly bounded operators.
Its proof uses both norm limits in \Cref{mod:means}: the Hilbert-space
limit controls one fourth-moment term, and the predual limit controls the
other.  No uniform operator norm bound is required for the initial
representatives of the spectral vectors.
Randomization, fourth-moment control, truncation, and selection of a
bounded witness occur together in the tracial spectral-gap argument of
Ioana and Vaes \cite[proof of the appendix theorem, pp.~129--130]{IoanaVaes2015}.
Related fourth-moment control gives equi-integrability in Marrakchi's
binormal-state implementation \cite[Section~2, Claim~2.2]{Marrakchi2025}.
We use the fourth-moment and polar-cutoff devices of Haagerup and Musat
\cite[Lemma~2.6 and Remark~2.8]{HaagerupMusat2007}, with moment bounds
supplied by modular orbit averages against the given state. The final
Fourier localization retains the shrinking spectral band.

\begin{theorem}[Bounded recovery]
\label{rec:bounded}
Let $M$ be a von Neumann algebra with a faithful normal state $\phi$ such
that $M_\phi=\C1$.  Use the standard-space notation and the fixed mean $m$
of \Cref{mod:average}.  Fix $T\in\B(H)$ and $s\in\R$.  Suppose that
$\delta_n>0$ tends to zero and that $h_n\in H$ are unit vectors with
$D$-spectral support in $[s-\delta_n,s+\delta_n]$.  If
\begin{equation}
\label{eq:recovery-positive}
 \limsup_{n\to\infty}m_t\norm{TU_th_n}^{2}>0,
\end{equation}
then there are a subsequence $(n_j)$, elements $v_j\in M$, and constants
$C_*<\infty$ and $\eta>0$ such that, for every $j$,
\[
 \norm{v_j}\le C_*,\qquad
 \norm{Tv_j\xi}\ge\eta,
\]
and $v_j\xi$ has $D$-spectral support in
$[s-4\delta_{n_j},s+4\delta_{n_j}]$.
\end{theorem}

\begin{proof}
We use the same family of Fourier filters to localize the initial bounded
approximations and to restore localization after polar truncation.
Choose $f_0\in C_c^\infty(\R)$ such that
\[
 0\le f_0\le1,\qquad f_0=1\text{ on }[-1,1],\qquad
 \supp f_0\subset[-2,2],
\]
and write $f_0(u)=\int_\R g_0(t)e^{itu}\,dt$, with $g_0$ a Schwartz
function.  For $\rho>0$, set
\[
 f_\rho(u)=f_0\bigl((u-s)/\rho\bigr),\qquad
 g_\rho(t)=\rho e^{-ist}g_0(\rho t).
\]
Then $f_\rho(u)=\int_\R g_\rho(t)e^{itu}\,dt$ and
\begin{equation}
\label{eq:recovery-kernel}
 \norm{g_\rho}_1=\norm{g_0}_1=:C_1
 \quad\text{for every }\rho>0.
\end{equation}
By \Cref{mod:localization}, the weak integral
\[
 \mathcal F_\rho(x)
   =\int_\R g_\rho(t)\sigma_t(x)\,dt\in M
\]
satisfies
\begin{equation}
\label{eq:recovery-filter}
 \mathcal F_\rho(x)\xi=f_\rho(D)x\xi,
 \qquad \norm{\mathcal F_\rho(x)}\le C_1\norm{x}.
\end{equation}
In particular, its vector has spectral support in
$[s-2\rho,s+2\rho]$.

\smallskip
\noindent\emph{Approximation by bounded representatives.}
By \Cref{eq:recovery-positive}, after passing to a subsequence and
reindexing, there is $\gamma>0$ such that
\[
 m_t\norm{TU_th_n}^{2}\ge4\gamma\qquad(n\ge1).
\]
Since $M\xi$ is dense in $H$, approximate $h_n$ by $x_n\xi$ and apply
$f_{\delta_n}(D)$.  This multiplier fixes $h_n$ and is a contraction.
Thus the vectors $\mathcal F_{\delta_n}(x_n)\xi$ approximate $h_n$ as
closely as desired.  Normalize them, obtaining $y_n\in M$ with
$\norm{y_n\xi}=1$ and spectral support in
$[s-d_n,s+d_n]$, where $d_n=2\delta_n$.  Choose the approximations so that
$\norm{y_n\xi-h_n}\to0$ and
\begin{equation}
\label{eq:recovery-defect}
 m_t\norm{TU_ty_n\xi}^{2}\ge2\gamma\qquad(n\ge1).
\end{equation}
This is possible because the averaged seminorm in
\Cref{mod:average} gives
\[
 \abs{\norm{t\mapsto TU_ty_n\xi}_{\mathrm{av}}
       -\norm{t\mapsto TU_th_n}_{\mathrm{av}}}
 \le\norm{T}\norm{y_n\xi-h_n}.
\]
Discard finitely many terms so that $d_n\le1$.  The modular identity
in \Cref{eq:modular-identities} gives
\begin{equation}
\label{eq:recovery-two-masses}
 \phi(y_n^*y_n)=1,\qquad
 \phi(y_ny_n^*)=\norm{e^{D/2}y_n\xi}^{2}\le e^{s+d_n}.
\end{equation}
There is no asserted uniform bound on $\norm{y_n}$.

\smallskip
\noindent\emph{Simultaneous defect and moment bounds.}
Fix $n$ for the moment, and write $y=y_n$.  With $\mu_k$ the uniform
probability measure on $[-k,k]$, define
\[
 \begin{aligned}
 A_k&=\int\sigma_t(y^*y)\,d\mu_k(t),&
 B_k&=\int\sigma_t(yy^*)\,d\mu_k(t),\\
 \psi(x)&=\phi(y^*xy),&
 \psi_k&=\int\psi\circ\sigma_{-t}\,d\mu_k(t).
 \end{aligned}
\]
We seek simultaneous control of $\phi(A_k^2)$ and $\psi_k(B_k)$, which
will bound the fourth moment after discretization. Here $\psi$ is a
normal state. State invariance gives the exact identity
\begin{equation}
\label{eq:recovery-functional}
 \psi_k(x)=
 \int\phi\bigl(\sigma_t(y)^*x\sigma_t(y)\bigr)\,d\mu_k(t).
\end{equation}
By the two ordinary norm limits in \Cref{mod:means},
\[
 A_k\xi\longrightarrow\xi,\qquad
 \norm{\psi_k-\phi}_{M_*}\longrightarrow0.
\]
Consequently $\phi(A_k^2)=\norm{A_k\xi}^{2}\to1$.  Although $B_k$
depends on $k$, it satisfies
\[
 \norm{B_k}\le\norm{y}^{2},\qquad
 \phi(B_k)=\phi(yy^*)\le e^{s+1},
\]
and hence
\begin{equation}
\label{eq:recovery-varying-test}
 \abs{\psi_k(B_k)-\phi(B_k)}
 \le\norm{\psi_k-\phi}_{M_*}\norm{y}^{2}\longrightarrow0.
\end{equation}
All these limits are taken with $n$ fixed.  For every sufficiently large
$k$ we may impose $\phi(A_k^2)\le2$ and
$\norm{\psi_k-\phi}_{M_*}\le\norm{y}^{-2}$ simultaneously.  It follows that
\begin{equation}
\label{eq:recovery-continuous-moments}
 \phi(A_k^2)+\psi_k(B_k)\le3+e^{s+1}.
\end{equation}
On the other hand, \Cref{eq:recovery-defect} and the definition of $m$
imply that the set of integers $k$ for which
\[
 \int\norm{T\sigma_t(y)\xi}^{2}\,d\mu_k(t)>\frac32\gamma
\]
belongs to the fixed free ultrafilter.  It therefore meets the cofinite
set on which \Cref{eq:recovery-continuous-moments} holds.  Choose one such
$k$, still for this fixed $n$.

We next replace this continuous average by a finite probability average.
Put $z_t=\sigma_t(y)$.  The three scalar functions
\[
 t\longmapsto\norm{Tz_t\xi}^{2},\qquad
 (t,u)\longmapsto\phi(z_t^*z_tz_u^*z_u),\qquad
 (t,u)\longmapsto\phi(z_t^*z_uz_u^*z_t)
\]
are continuous on the relevant compact interval or square.  Indeed,
$t\mapsto z_t$ is bounded and strong-star continuous, and products on
bounded sets preserve this continuity.  A common sufficiently fine
partition of $[-k,k]$, using its normalized interval lengths as weights,
therefore approximates all three integrals.  The product weights
approximate the double integrals defining $\phi(A_k^2)$ and
$\psi_k(B_k)$.  Thus there are finitely many times $t_i$ and weights
$p_i\ge0$, with $\sum_i p_i=1$, such that, writing
\[
 z_i=\sigma_{t_i}(y),\qquad
 A=\sum_i p_i z_i^*z_i,\qquad
 B=\sum_i p_i z_iz_i^*,
\]
we have
\begin{equation}
\label{eq:recovery-finite}
 \sum_i p_i\norm{Tz_i\xi}^{2}\ge\gamma,
 \qquad
 \phi(A^2)+\sum_i p_i\phi(z_i^*Bz_i)\le C,
 \qquad C:=4+e^{s+1}.
\end{equation}
The finite averages may depend on $n$; $\gamma$ and $C$ do not.  In
particular, no estimate on the rate in
\Cref{eq:recovery-varying-test} uniform in $n$ has been used.

\smallskip
\noindent\emph{Randomization and truncation.}
For this finite average, let $\epsilon_i$ be independent random variables
uniformly distributed on the complex unit circle, and set
\[
 z=\sum_i\sqrt{p_i}\,\epsilon_i z_i.
\]
The vector $z\xi$ is supported in $[s-d_n,s+d_n]$.  Independence and
vanishing first moments give
\[
 \mathbb E\norm{Tz\xi}^{2}
   =\sum_i p_i\norm{Tz_i\xi}^{2}\ge\gamma.
\]
The following expansion is the Steinhaus variant of the fourth-moment identity in
\cite[equation~(2.10) and Remark~2.8]{HaagerupMusat2007}, with
coefficients $\sqrt{p_i}z_i$. We include the expansion to keep track of
the order of the noncommuting factors. The two surviving pairings give
\begin{equation}
\label{eq:recovery-fourth}
 \mathbb E[(z^*z)^2]
  =A^2+\sum_i p_i z_i^*Bz_i
       -\sum_i p_i^2(z_i^*z_i)^2.
\end{equation}
The last sum subtracts the overlap of the pairings, where all four
indices agree.  Its summands are positive, so
\Cref{eq:recovery-finite} implies
\begin{equation}
\label{eq:recovery-fourth-bound}
 \mathbb E\phi((z^*z)^2)\le C.
\end{equation}

We next use the functional-calculus truncation underlying
\cite[Lemma~2.6, equations~(2.15)--(2.20)]{HaagerupMusat2007}.
For $K>0$, write $z=u\abs{z}$ in its polar decomposition and put
$z^{[K]}=u\min(\abs{z},K)$.  Then $\norm{z^{[K]}}\le K$ and
\begin{equation}
\label{eq:recovery-truncation}
 \begin{aligned}
 \mathbb E\norm{(z-z^{[K]})\xi}^{2}
 &=\mathbb E\phi\bigl((\abs{z}-K)_+^2\bigr)\\
 &\le K^{-2}\mathbb E\phi(\abs{z}^{4})
 \le CK^{-2}.
 \end{aligned}
\end{equation}
This controls the vector error on the side required for the bounded
operator $T$.  No estimate for the adjoint truncation error is needed.

\smallskip
\noindent\emph{Restoring the spectral band and selecting realizations.}
Define the random element $v=\mathcal F_{d_n}(z^{[K]})\in M$.
By \Cref{eq:recovery-filter}, every realization satisfies
\begin{equation}
\label{eq:recovery-bounded-output}
 \norm{v}\le C_1K,
 \qquad v\xi=f_{d_n}(D)z^{[K]}\xi,
\end{equation}
and $v\xi$ is supported in $[s-2d_n,s+2d_n]$.  Since $f_{d_n}(D)$ fixes
$z\xi$ and has norm at most one, \Cref{eq:recovery-truncation} gives
\[
 \mathbb E\norm{v\xi-z\xi}^{2}\le CK^{-2}.
\]
The triangle inequality in the Hilbert space of square-integrable random
$H$-vectors now yields
\begin{equation}
\label{eq:recovery-fixed-cutoff}
 \bigl(\mathbb E\norm{Tv\xi}^{2}\bigr)^{1/2}
 \ge\sqrt\gamma-\frac{\norm{T}\sqrt C}{K}.
\end{equation}
Fix once and for all
\[
 K\ge2\norm{T}\sqrt{C/\gamma}.
\]
The hypothesis implies $T\ne0$, so $K$ can be chosen positive.  This
choice is independent of $n$ and ensures that the left-hand side of
\Cref{eq:recovery-fixed-cutoff} is at least $\sqrt\gamma/2$.

To select a realization, note that $z$ depends norm-continuously on the
finite phase torus and $z^{[K]}=zq_K(z^*z)$, where
\[
 q_K(r)=\min\{1,K/\sqrt r\}\quad(r>0),\qquad q_K(0)=1.
\]
Since $q_K$ is continuous on $[0,\infty)$, continuous functional calculus
makes the cutoff norm-continuous. The filter $\mathcal F_{d_n}$ is bounded
and linear, so $\norm{Tv\xi}^2$ is continuous on the compact phase torus.
It therefore attains a maximum at least its expectation.
Choose such a realization $v_n$ for each $n$. Then
\[
 \norm{v_n}\le C_1K,
 \qquad \norm{Tv_n\xi}\ge\frac{\sqrt\gamma}{2},
\]
and its spectral support lies in
$[s-2d_n,s+2d_n]=[s-4\delta_n,s+4\delta_n]$.
Restoring the subsequence notation proves \Cref{rec:bounded} with
$C_*=C_1K$ and $\eta=\sqrt\gamma/2$.
\end{proof}

\Cref{fig:bounded-recovery} summarizes the construction. The averaging
time depends on the initial representative, whereas the moment bound,
cutoff threshold, and Fourier-kernel norm are independent of the width
of its spectral band.
\begin{figure}[tbp]
\centering
\begin{tikzpicture}[
  stage/.style={draw=black!55,rounded corners=2pt,fill=blue!3,
    text width=3.75cm,minimum height=1.85cm,align=center,
    inner sep=5pt,font=\small},
  flow/.style={-{Stealth[length=2mm]},semithick,draw=black!70}]
  \node[stage] (input) at (0,0) {
    \textbf{Spectral vectors}\\[3pt]
    $\norm{h_n}=1$, radius $\delta_n$\\
    $m_t\norm{TU_th_n}^2\ge4\gamma$};
  \node[stage] (approx) at (4.65,0) {
    \textbf{Operator vectors}\\[3pt]
    $y_n\xi$, radius $2\delta_n$\\
    $\norm{y_n}$ may grow with $n$};
  \node[stage] (random) at (9.3,0) {
    \textbf{Average and randomize}\\[3pt]
    $z=\sum_i\sqrt{p_i}\epsilon_i z_i$\\
    $\mathbb E\phi(|z|^4)\le C$};
  \node[stage] (cutoff) at (9.3,-2.35) {
    \textbf{One polar cutoff}\\[3pt]
    $\norm{z^{[K]}}\le K$\\
    mean squared error\\
    $\le C/K^2$};
  \node[stage] (filter) at (4.65,-2.35) {
    \textbf{Restore the spectral band}\\[3pt]
    $v=\mathcal F_{2\delta_n}(z^{[K]})$\\
    $\norm{v}\le C_1K$, radius $4\delta_n$};
  \node[stage] (output) at (0,-2.35) {
    \textbf{Select a realization}\\[3pt]
    $\norm{v_n}\le C_1K$\\
    $\norm{Tv_n\xi}\ge\sqrt\gamma/2$};
  \draw[flow] (input) -- (approx);
  \draw[flow] (approx) -- (random);
  \draw[flow] (random) -- (cutoff);
  \draw[flow] (cutoff) -- (filter);
  \draw[flow] (filter) -- (output);
\end{tikzpicture}
\caption{The construction in \Cref{rec:bounded}, after passing to the
subsequence with a positive averaged squared norm. Every band is centered at the
fixed modular energy $s$. The polar cutoff need not preserve that band;
the next Fourier filter restores it with radius $4\delta_n$. The constants
$C$, $K$, and $C_1$ are independent of $n$. The cutoff error is averaged
over the circle phases.}
\label{fig:bounded-recovery}
\end{figure}

